%!TEX root = TQFTbook.tex
%%%%%%%%%%%%%%%%%%%%%%%%%%%%%%%%%%%%%%%%%%%%%%%%%%%%%%%%%%
\chapter{Physical motivation}\label{c:physical-motivation}
In this chapter we provide a motivation for the mathematical study of
topological quantum field theories (TQFTs). These theories  arise as a special 
case of more general Quantum Field Theory (QFT), so  we first look at the physics
definition of QFT.  This definition is notoriously non-rigorous, and as of now,
there is no mathematically rigorous definition of QFT in full generality.
Despite this, one can use the non-rigorous physical definition to try to
axiomatize the bare minimum requirements that any sensible QFT should fulfill.

We start with a heuristic introduction to classical and quantum field theories.
Since our purpose is just to motivate the connection to TQFT, we suppress here
many technical details to keep the exposition simple. Interested readers can
find details in numerous textbooks on QFT, such as physics classic
\ocite{peskin-schroeder} or a recent book \ocite{etingof-qft},
written for mathematicians.

This chapter is used as a motivation only; no results in the remaining
chapters rely on the material of this chapter.


%%%%%%%%%%%%%%
\section{Classical and quantum mechanics}\label{s:mechanics}

In classical mechanics, at any time $t$, a point particle is described by its
position $x(t) \in N$ and its velocity $\dot x(t) \in T_{x(t)}N$. Here $N$ is
the configuration space which, for a point particle, is the same as physical
space, and  $T_{x(t)}N$ is the tangent space to $N$ at point $x(t)$.
(For a more general mechanical system, $N$ can be some other space of
configurations.)
The trajectory of a particle is a map from the
time interval $[t_0,t_1]$ to $N$:
\begin{equation}
	\ga\colon[t_0,t_1] \rightarrow N, \nonumber
\end{equation}
We are using the position variable $x$ and the trajectory $\ga$
interchangeably. They are both continuous maps from the time-interval to the
configuration space, but the notation $\ga$ is better suited for quantum
mechanics where we will be considering all such maps.

The equations of motion are usually derived using the Lagrangian formalism.
Namely, for each trajectory $\ga(t)$ we define its  action functional:
\begin{equation}
	S[\ga]= \int L(x(t),\dot x(t)) dt \nonumber
\end{equation}
where $L(x, \dot x)$ is a smooth function on $TN$ called the Lagrangian.
 The precise form of
Lagrangian function $L$ depends on the physics governing the motion of the
particle.

The  least action principle (more appropriately called the stationary action
principle) states that the actual trajectory that the physical particle takes
is one for which the functional variation of the action vanishes (in the class
of all trajectories with fixed starting and ending position):
\begin{equation}\label{e:stationary-path}
	\delta S[ x_{ph}(t)] = 0 , x_{ph}(t_0) = x_0, x_{ph}(t_1) = x_1
\end{equation}
This condition leads to so-called Euler--Lagrange equations, which are the
equations of motion in classical mechanics.


In quantum mechanics, the state of the particle is described by the wave
function $\psi$, a vector  in some Hilbert space $\calH_N$; for a single
particle moving in configuration space $N$, $\calH_N$ is the space of square
integrable functions on $N$:
\begin{equation*}
	\psi \in \calH_N \cong L^2(N).
\end{equation*}

The evolution of a quantum-mechanical system from time $t_0$
to $t_1$ is described  by the evolution operator:
\begin{equation*}
	U_{t_0,t_1}\colon \calH_N \to\calH_N.
\end{equation*}

Traditionally, this operator is defined using the Schr\"odinger
equation. However, there is an alternative approach, based on Feynman's path
integral. Namely, let us  express the evolution operator using an integral kernel:
\begin{equation}
	(U_{t_0,t_1} \psi)(x_1) = \int_{N} K(x_1,x_0) \psi(x_0) dx_0
\end{equation}

The Feynman path integral prescription states that this integral 
kernel can be written as:
\begin{equation}\label{e:feynman1}
	K(x_1,x_0) = \int_{\substack{\ga(t_0)=x_0 \\
  \ga(t_1)=x_1}} e^{ \frac{i}{\hbar} S[\ga] } D\ga
\end{equation}
where $\hbar$ is the Planck constant.

The physical essence of this formula is that all possible paths taking the
particle from $x_0$ at $t_0$ to $x_1$ at $t_1$ contribute to its evolution.
This prescription of integration over all trajectories is sometimes called
``sum over histories''.

Formula \eqref{e:feynman1} expresses the evolution operator in terms of
integral over the infinite-dimensional space of paths. Making sense of such
integrals is non-trivial; however, heuristically one can argue that  in the
limit $\hbar \to 0$, the integral in \eqref{e:feynman1} is dominated by the
contribution from the neighborhood of critical points of the action. For
finite-dimensional integrals of this form, method of stationary phase allows
one to write asymptotic expansion in powers of $\hbar$; using that as
inspiration, one can also write similar formulas for asymptotic expansion of
integral \eqref{e:feynman1}.  The lowest order term in this expansion
corresponds to contribution from the classical path defined in
\eqref{e:stationary-path}.

In Table \ref{tab:mechanics}, we summarize the main structures of classical
mechanics and the analogous quantities in quantum mechanics. In the next
section, we will use this prescription of sum over histories to define time
evolution in a quantum field theory.

\begin{table}[ht]
  \centering
  \begin{tabular}{|r|l|l|}
  	 \hline
    &\textbf{Classical mechanics} & \textbf{Quantum mechanics} \\[5pt]
    \hline
    Position
    &$x\in N$ (configuration space)
    & $\psi \in \calH_N \cong L^2(N)$ \\[5pt]
    Trajectory
    &$\gamma\colon [t_0, t_1]\to N$
    &$U_{t_0,t_1}\colon \calH_N\to \calH_N$ \\[5pt]
    Action
      &\multicolumn{2}{c|} {$S[\gamma]=\int L(\gamma,\dot\ga) dt$}
      %& $K(x_1,x_0) = \int e^{\frac{i}{\hbar} S[\gamma]} D\ga$ 
      \\[5pt]
    Laws of motion
    &\parbox[t]{5cm}{Stationary action principle\par
    	\quad $\delta S[\ga]=0$ (fixed endpoints)
    }
    & \parbox[t]{5.5cm}{
    	Sum over histories \par
    	$(U\psi)(x_1) = \int_N K(x_1, x_0) \psi(x_0) dx_0$\par
       $K(x_1,x_0) = \int e^{\frac{i}{\hbar} S[\gamma]} D\ga$
    }
      \\
    \hline
  \end{tabular}
\medskip
  \caption{Summary of the main structures in classical and quantum mechanics}
  \label{tab:mechanics}
\end{table}


%%%%%%%%%%%%%%
\section{Classical and quantum field theory}\label{s:field-theory}
Let us now move from mechanics to field theory.  In field theory,
the position variable $x\in N$ is replaced by a field $\ph$,
which is typically either a function on $N$ (with values in some target space
$V$),  a section of a vector bundle, or a connection. For simplicity, let us
assume that the field $\ph$ (at a given time) takes values in $V$:
\begin{equation*}
	\ph\colon N \rightarrow V.
\end{equation*}
As is common in the physics literature, we will assume that the field $\ph$
is smooth enough for all purposes. We denote by $\calF_N$ the space of all
fields (at fixed moment of time); in the simplest example above,
$\calF_N=C^\infty (N,V)$. This space is the analog of configuration space
$N$ in classical mechanics.

The analog of a trajectory is now a field $\ph(x,t)$ on spacetime $M = N
\times [t_0,t_1]$. If the space $N$ is $(n-1)$-dimensional, then
spacetime $M$ is $n$-dimensional; the corresponding field theory is referred
to as $n$-dimensional field theory.


 Note that the space $N$ plays a different role here than in the mechanics. In
 a field theory, $N$ is the domain of field $\ph$, whereas in mechanics, the
 configuration space $N$, was the target of ``position field'' $x$. From this
 perspective, we can think of the discussion in the previous section as a 0+1
 dimensional field theory of the position field $x\colon \text{point}
 \to N$.

As in mechanics, equations of motion come from the stationary action principle.
Namely, for a field $\ph$ defined on spacetime $M=N\times[t_0, t_1]$, we define
the action functional by
\begin{equation}
	S[\ph] = \int_{t_0}^{t_1} L(\ph, \del_{\mu} \ph) dt = \int_{t_0}^{t_1} \left(
	\int_N \calL(\ph, \del_{\mu} \ph) d^{n-1} x \right) dt
\end{equation}
$\calL$ is called the Lagrangian density; it depends on the values of field $\ph$ 
at point $(x,t)$, its first derivatives $\del_\mu \ph$, and on metric $g$ on spacetime. 

The classical equations of motion again follow from the stationary action
principle: $\delta S[\ph]=0$. This is equivalent to a system of PDE's on
$\ph$, called the \emph{Euler--Lagrange equations}.
%\begin{equation}
%	\del_{\mu} \left( \frac{\del \calL}{\del (\del_{\mu} \ph)} \right) -
%\frac{\del \calL}{\del \ph} = 0
%\end{equation}

%This easily generalizes to a Lagrangian that depends on more than one field.
%The equations coming from $\delta S[\ph]=0$ are generally called the
%Euler-Lagrange equations of motion.

Similar to how we moved from classical to quantum mechanics, it is natural to
expect that in quantum field theory, the state of the system is described by a
functional $\psi \in L^2(\calF_N)$. The space $\calF_N$ is
infinite-dimensional, and there is no canonical Lebesgue measure that we can
define for it, so exact meaning of $L^2$ is not well-defined, but for the
moment, let us ignore that.


The time evolution operator $U_{t_0,t_1}\colon \calH_N \rightarrow \calH_N$ 
can again be defined using the path integral.
Given a state $\psi$ at time $t_0$, the final state $U\psi$ at time $t_1$
can be written as:

\begin{equation}\label{e:time-evolution}
  \begin{aligned}
	(U \psi)[\ph_1] &= \int_{\calF_N} K(\ph_1,\ph_0) \psi[\ph_0] D \ph_0\\
	K(\ph_1,\ph_0) &= \int_{\substack{\Phi|_{t_0}
  =\ph_0 \\ \Phi|_{t_1}=\ph_1}} e^{\frac{i}{\hbar} S[\Phi]} D\Phi.
  \end{aligned}
\end{equation}
Here $\Phi=\Phi(x,t)$ is a field on spacetime $M=N\times[t_0,t_1]$ with 
given values at $t=t_0$, $t=t_1$; it is an analog of  path $\ga$ in  mechanics, and 
$D\Phi$ is some (not yet defined) measure on the space of fields. 


Note that the integral in \eqref{e:time-evolution} does not explicitly use 
the decomposition of spacetime $M$ as a product of space $N$ and time interval 
$[t_0, t_1]$: it only uses the metric on $M$ (which is used in the Lagrangian density)
and the fact that  $M$ has two boundary components, $N_0=N\times\{t_0\}$ and 
$N_1 = N\times \{t_1\}$. The same integral makes perfect sense if we allow 
spacetime $M$ to have a different topology: for any $n$-dimensional manifold $M$ 
whose boundary is presented as $\del M=\ov{N_0}\sqcup N_1$ (where bar stands for 
orientation reversal), formula \eqref{e:time-evolution} defines an operator 
$U\colon \calH_{N_0}\to \calH_{N_1}$. 



The above formulas freely use integrals over infinite-dimensional spaces of
paths/fields, which are very hard to define rigorously; there are many
continuum quantum field theories for which a complete rigorous construction of
the path integral is not currently available. Usually the best one can do is to
write the asymptotic formulas, expanding the operator $U$ or the kernel $K$ as
a series in powers of $\hbar$, and even that requires a lot of work. 

However, if one is willing to ignore these problems for now, then we can
use path integral formula \eqref{e:time-evolution} for showing various
properties of QFT. For example,
it is easy to show the composition property
$$
U_{t_0,t_2}= U_{t_1,t_2} U_{t_0,t_1}
$$
Indeed, since a field $\Phi\in \calF_{N\times [t_0,t_2]}$ is the same as a 
pair of fields $\Phi_1\in \calF_{N\times [t_0,t_1]}$, 
$\Phi_2\in \calF_{N\times [t_1,t_2]}$ whose values at the common boundary 
$N\times\{t_1\}$ coincide, we get 
 
\begin{equation}\label{e:gluing}
  K_M(\ph_2,\ph_0) = \int_{\calF_N}
  K_{M_2}( \ph_2,\ph_1) K_{M_1}(\ph_1, \ph_0) D \ph_1
\end{equation}
where $M_1=N\times [t_0,t_1]$,  $M_2=N\times [t_1,t_2]$. 


More generally, the same (non-rigorous) argument shows that for more general 
spacetimes $M$, we can choose an arbitrary codimension one slice $N$, and split the path integral into two components of $M$ that share $N$ as a common boundary.



In the next section, we will try to bypass these rigor-related problems by
approaching QFT in a more axiomatic manner. The main structures appearing
in classical and quantum field theories are summarized in Table
\ref{tab:qft}.

\begin{table}[h]
	\centering
	\begin{tabular}{|r|l|l|}
		\hline
		&\textbf{Classical field theory} & \textbf{QFT} \\[5pt]
		\hline
		Field/State &$\ph \in \calF_N$
		& $\psi \in \calH_N \cong L^2(\calF_N)$ \\[5pt]
		Trajectory
		&$\ph\in \calF_{M}$, $M=N\times[t_0, t_1]$
		&$U_{t_0,t_1}\colon \calH_N\to \calH_N$ \\[5pt]
		Action
		& \multicolumn{2}{c|}{$S[\ph]=\int_M \calL(\ph, \del_\mu\ph) d^{n}x$}
		%&$K(\ph_1,\ph_0) = \int e^{\frac{i}{\hbar} S[\Phi]} D\Phi$ 
    \\[5pt]
		Laws of motion
		&\parbox[t]{5cm}{Euler-Lagrange equations\par
			\quad $\delta S[\ph(t)]=0$
		}
		& \parbox[t]{6cm}{
			$(U \psi)[\ph_1] = \int_{\calF_N} K(\ph_1,\ph_0) \psi[\ph_0] D \ph_0$\par
      $K(\ph_1,\ph_0) = \int e^{\frac{i}{\hbar} S[\Phi]} D\Phi$
		}\\
		\hline
	\end{tabular}
	\caption{A schematic comparison of main structures in classical 
    and quantum field theory. See main text for details.}
	\label{tab:qft}
\end{table}

%%%%%%%%%%%%%%%%%%%%%%%%%%%
\section{Axioms of quantum field theory}\label{s:axioms-field-theory}

Let us ignore the problems involved in the above definition of QFT, and let us
only focus on the formal properties that any sensible quantum field theory
should satisfy, focusing on a few
axioms that are necessary to transition to TQFTs. 

Taking motivation from the path integral formulation of QFT outlined in the previous section, 
we define a QFT as the following structure:
\begin{itemize}
  \item For every $(n-1)$-dimensional manifold $N$, we have a \emph{Hilbert space} $\calH_N$.
  \item For every $n$-dimensional manifold $M$ together with decomposition 
    $\del M=\ov{N_0}\sqcup N_1$, we have the evolution operator $Z_M\colon
   \calH_{N_0}\to \calH_{N_1}$, given by the path integral formula
   \eqref{e:time-evolution}. (Previously, we used $U$ for the time
  evolution; we switch the notation to $Z$ which is more common in QFT).
\end{itemize}

These two pieces of data should satisfy certain conditions for consistency:

\begin{enumerate}
  	\item {\bf Disjoint union axiom}: Assume that $N$ is disconnected: 
      $N=N_1\sqcup N_2$.  Then for the space of fields, we have 
      $\calF_{N}=\calF_{N_1}\times \calF_{N_2}$. Using the heuristic definition 
      $\calH_N=L^2(\calF_N)$, we therefore require 
    \begin{equation}\label{e:qft_axiom1}
      \calH_{N_1\sqcup N_2} \cong \calH_{N_1}\otimes \calH_{N_2}
    \end{equation}

  \item {\bf Gluing axiom}: So defined evolution operators satisfy the gluing axiom, 
    which is a generalization of \eqref{e:gluing}.

    Consider a manifold $M$ that can be cut into two manifolds $M_1$ and $M_2$ along a codimension one submanifold $N$, see \firef{f:gluing} for an
    illustration. If $\partial M = \overline{N_0} \sqcup N_1$, then $\del M_{1} = \overline{N_0} \sqcup N$, and $\del M_{2} = \overline{N} \sqcup N_1$. 

    Gluing axiom says that operator $Z_{M}$ is the same as composing operators $ Z_{M_1}$ and $Z_{M_2}$:
    \begin{equation}\label{e:gluing-axiom}
    	Z_{M} = Z_{M_2} \circ Z_{M_1}
    \end{equation}
	In the path integral picture, this is a natural expectation. Instead of doing path integral over entire manifold $M$, we can do the path integral in two ( or more ) steps.  Evolving a state $\psi$ from $\mathcal{H}_{N_0}$ to $\mathcal{H}_{N}$ with operator $Z_{M_1}$, and then evolving the state $Z_{M_1} \psi$ from $\mathcal{H}_{N}$ to $\mathcal{H}_{M_1}$ with operator $Z_{M_1}$ produces the state $( Z_{M_2} \circ Z_{M_1} ) \psi$. We expect this final state to equal $Z_M \psi$ for every initial state $\psi$.
	
      Note that the manifolds don't need to be cylinders. If we consider the case where $M_1=N\times [t_0,t_1]$,  $M_2=N\times [t_1,t_2]$, then we recover \eqref{e:gluing} from the previous section.
    \begin{figure}[ht]
    	\centering

    	\begin{tikzpicture}[scale=0.25]
    		\begin{scope}[xshift=0cm]
    			% first diagram code
    			\node (31) at (-9.5, -1.75) {};
    			\node (33) at (3.25, 0.75) {};
    			\node (34) at (-7.5, 1) {};
    			\node (35) at (-4.25, 1.25) {};
    			\node (36) at (-7, 0.75) {};
    			\node (37) at (-4.75, 1) {};
    			\node (38) at (5.5, 0.25) {};
    			\node (41) at (5.5, -2.75) {};
    			\node (42) at (2.75, -2) {};
    			\node (43) at (-2, -1.75) {};
    			\node (45) at (-5.75, 1.5) {};
    			\node (46) at (5.25, 5) {};
    			\node (47) at (5.25, 2.5) {};
    			\node (48) at (3.5, 2.25) {};
    			\node (49) at (1, 3.5) {};
    			\node (50) at (-2.25, 3.5) {};
    			\node (51) at (-9.5, 3) {};
    			\node (52) at (-5.5, 4.5) {};
    			\node (53) at (-0.75, 2.83) {};
    			\node (54) at (-0.75, -1.22) {};
    			\draw [bend right, looseness=1.25] (34.center) to (35.center);
    			\draw [bend left=60, looseness=1.25] (36.center) to (37.center);
    			\draw [bend right, looseness=0.75] (33.center) to (38.center);
    			\draw [bend left=15, looseness=1.25] (38.center) to (41.center);
    			\draw [bend right, looseness=1.25] (31.center) to (43.center);
    			\draw [bend left=45, looseness=0.75] (43.center) to (42.center);
    			\draw [bend right, looseness=0.75] (42.center) to (41.center);
    			\draw [bend right=15, looseness=1.25] (38.center) to (41.center);
    			\draw [bend left=15, looseness=1.25] (46.center) to (47.center);
    			\draw [bend right=15, looseness=1.25] (46.center) to (47.center);
    			\draw [bend left=330, looseness=0.75] (47.center) to (48.center);
    			\draw [bend right=45] (48.center) to (33.center);
    			\draw [bend left=345, looseness=1.25] (46.center) to (49.center);
    			\draw [bend right=315] (49.center) to (50.center);
    			\draw [bend left=15] (52.center) to (50.center);
    			\draw [bend left=15, looseness=0.75] (51.center) to (52.center);
    			\draw [bend right=45, looseness=0.50] (51.center) to (31.center);
    			\draw [bend left=60, looseness=0.50] (51.center) to (31.center);
    			\node[left] at ($(45.center)+(+5.00,-0.00)$) {$M$};
    			\node[left] at ($(45.center)+(+12.50,-0.50)$) {$N_0$};
    			\node[left] at ($(45.center)+(-4.00,-0.50)$) {$N_1$};
    			
    			
    		\end{scope}
    		
    		\begin{scope}[xshift=20cm]
    			% second diagram code
    			\node (31) at (-9.5, -1.75) {};
    			\node (33) at (3.25, 0.75) {};
    			\node (34) at (-7.5, 1) {};
    			\node (35) at (-4.25, 1.25) {};
    			\node (36) at (-7, 0.75) {};
    			\node (37) at (-4.75, 1) {};
    			\node (38) at (5.5, 0.25) {};
    			\node (41) at (5.5, -2.75) {};
    			\node (42) at (2.75, -2) {};
    			\node (43) at (-2, -1.75) {};
    			\node (45) at (-5.75, 1.5) {};
    			\node (46) at (5.25, 5) {};
    			\node (47) at (5.25, 2.5) {};
    			\node (48) at (3.5, 2.25) {};
    			\node (49) at (1, 3.5) {};
    			\node (50) at (-2.25, 3.5) {};
    			\node (51) at (-9.5, 3) {};
    			\node (52) at (-5.5, 4.5) {};
    			\node (53) at (-0.75, 2.83) {};
    			\node (54) at (-0.75, -1.22) {};
    			\draw [bend right, looseness=1.25] (34.center) to (35.center);
    			\draw [bend left=60, looseness=1.25] (36.center) to (37.center);
    			\draw [bend right, looseness=0.75] (33.center) to (38.center);
    			\draw [bend left=15, looseness=1.25] (38.center) to (41.center);
    			\draw [bend right, looseness=1.25] (31.center) to (43.center);
    			\draw [bend left=45, looseness=0.75] (43.center) to (42.center);
    			\draw [bend right, looseness=0.75] (42.center) to (41.center);
    			\draw [bend right=15, looseness=1.25] (38.center) to (41.center);
    			\draw [bend left=15, looseness=1.25] (46.center) to (47.center);
    			\draw [bend right=15, looseness=1.25] (46.center) to (47.center);
    			\draw [bend left=330, looseness=0.75] (47.center) to (48.center);
    			\draw [bend right=45] (48.center) to (33.center);
    			\draw [bend left=345, looseness=1.25] (46.center) to (49.center);
    			\draw [bend right=315] (49.center) to (50.center);
    			\draw [bend left=15] (52.center) to (50.center);
    			\draw [bend left=15, looseness=0.75] (51.center) to (52.center);
    			\draw [bend right=45, looseness=0.50] (51.center) to (31.center);
    			\draw [bend left=60, looseness=0.50] (51.center) to (31.center);
    			\draw [bend right=15] (53.center) to (54.center);
    			\draw [dashed, bend left=15] (53.center) to (54.center);
    			\node[left] at ($(45.center)+(9.00,0.80)$) {$M_1$};
    			\node[right] at ($(45.center)+(+0.00,0.90)$) {$M_2$};
    			\node[left] at ($(45.center)+(+5.00,-0.50)$) {$N$};
    			\node[left] at ($(45.center)+(+12.50,-0.50)$) {$N_0$};
    			\node[left] at ($(45.center)+(-4.00,-0.50)$) {$N_1$};
    			
    		\end{scope}
    		
    	\end{tikzpicture}

    	\caption{Cutting of a manifold $M$ into two manifolds $M_1$ and $M_2$, along a codimension one submanifold $N$. Gluing axiom asserts that $Z_M=Z_{M_2} \circ Z_{M_1}$. }
    	\label{f:gluing}
    \end{figure}

    Gluing axiom is a very strong axiom. It allows us to compute the path integral of a manifold from smaller pieces. The gluing axiom tells us how the (path integrals of) smaller pieces are to be glued along their boundaries in order to reconstruct the path integral of the original manifold. The underlying assumption that we can compute path integral of smaller pieces individually makes the locality of a theory manifest.



Note that for an empty space $\varnothing$, it is natural to expect
$\calH_\varnothing=\C$. This easily follows from the multiplication axiom since $\varnothing$ is the unit of disjoint union operation : $ \calH_{N} \cong \calH_{N \sqcup \varnothing} \cong \calH_{N}\otimes \calH_{\varnothing}  $. Thus, for a closed
spacetime (i.e., $\del M=\varnothing$), we get an operator
\begin{equation}
	Z_M\colon \C\to \C.
\end{equation}
In other words, $Z_M\in \C$ is a number; in physics it is usually called the 
\emph{partition function}\index{partition function} and denoted by $Z(M)$.


\end{enumerate}

\section{Topological QFT}\label{s:tqft}

So far we have not specified what we mean by ``manifold''. Usually in physics
the spacetime is an oriented pseudo-Riemannian manifold, with metric of
signature $(1, n-1)$. However, we will be interested in a very special case,
namely when all the structures above do not use the metric on the space; thus,
we only require $M$ to be an oriented smooth manifold. Such theories are called
\emph{topological}. The most famous example of such a theory is Chern--Simons
theory (see \ocite{witten-jones}), in which the spacetime $M$ must be
3-dimensional, and fields are $\SU(N)$ connections on $M$. After choosing a
trivialization of the bundle, these fields can be described by one-forms,
usually denoted by $A$, with values in the Lie algebra $\su(N)$. The action is
then given by
\begin{equation}
  S[A] = \frac{k}{4 \pi} \int_M \tr(A\wedge dA + \tfrac{2}{3}A\wedge A\wedge A)
\end{equation}

The `level' $k$ is a fixed integer parameter of the theory. Although the
numerical value of the Chern--Simons action may change by an integer multiple of
$2\pi k$ under a change of trivialization, the exponentiated action $e^{i
S[A]}$ is unchanged. The important thing to note here is that the path integrals
defined using this action do not depend on the choice of the metric. So the
partition functions assigned by this QFT are invariant under
orientation-preserving diffeomorphisms. (It turns out that there are other
observables in this theory besides the partition functions, namely the Wilson
loops, but our preliminary definition of a TQFT is agnostic of these
observables.) We can axiomatize the invariance under diffeomorphisms as a
defining property of a topological theory. Together with the axioms of a QFT,
we get the following preliminary definition:

\begin{predefinition}\label{d:ptqft}
  An $n$-dimensional topological field theory (TQFT) is the following
  collection of data
  \begin{itemize}
    \item For every closed $(n-1)$-dimensional oriented smooth manifold $N$, a
      vector space $Z_N$
    \item For every $n$-dimensional oriented manifold $M$ with boundary
      $\del M=\ov{N_0}\sqcup N_1$, a linear operator $Z_M\colon Z_{N_0}\to Z_{N_1}$,
  \end{itemize}

  such that

  \begin{itemize}
  	\item Orientation preserving diffeomorphisms preserve this data
  	\item This data satisfies the disjoint union axiom and the gluing axiom defined above.
  \end{itemize}

\end{predefinition}

There is a lot more detail to be filled in, but that requires some preparation. 
We will provide a more precise definition, due to Atiyah \ocite{atiyah}, later (see \deref{d:tqft}).

Remarkably, topological theories are much easier to define rigorously than general QFTs. 
For example, as we show later, in a TQFT all Hilbert spaces assigned to $(n-1)$-dimensional
manifolds are automatically finite-dimensional, thus avoiding all analytical difficulties with 
convergence. And  instead of defining a TQFT by a path integral or other
global construction (which can be very hard to do), we can try to construct
TQFTs by defining $Z(M)$ for ``elementary blocks'' and gluing every manifold
from them. For example, for $n=2$, any manifold can be obtained by gluing copies of 
the  disk, cylinder, and pair of pants.
\begin{center}
  \begin{tikzpicture}
    \pic at (0, -0.5) {cup};
    \pic at (1.5,0) {cylinder};
    \pic at (3,0) {pants};
  \end{tikzpicture}
\end{center}

Of course, we need to have some compatibility conditions to ensure that
different ways of gluing the manifold from pieces give the same result;
essentially, we are trying to define manifolds by ``generators and
relations''. For $n=2$ it is not too hard (we will do so in
\chref{c:2d-classification}); for larger $n$, it is harder --- e.g., for
$n=3$ we can consider so-called Heegaard decomposition (which already
requires infinitely many blocks, namely all handlebodies of arbitrary
genus) or use surgery along knots. The difficulty lies in how fast the data
grows with higher dimensions. Ideally, we would want the number of
generators/relations to be as small as possible in a TQFT, or at least
finite. One such notion is that of an extended TQFT.

%%%%%%%%%%%%%%%%%%%%%%%%%%%%%%%%%%%%%%%%%%
\section{Toward extended TQFT}\label{s:xtqft}
So far we talked about cutting a manifold into smaller manifolds with
boundaries. By assigning appropriate data to manifolds with boundaries, we
could glue them back and assign a topological invariant to all possible closed
$n$-manifolds. But why stop at manifolds with boundary? If we allow manifolds
with corners of any codimension, the set of ``elementary blocks'' is greatly
simplified: any manifold can be glued together from just one kind of local
block, namely an $n$-simplex. Such theories are called fully extended (or
sometimes local) TQFTs; see \ocite{lurie-cobordism}.

But the price we pay is algebraic complexity. In our discussion so far, a TQFT
assigned a number to a closed $n$-manifold, and a vector space to a
codimension-one manifold. If we allow pieces with codimension 2 manifolds, what
kind of data should be assigned to them? And then what about the next level?

One way to answer is that such higher codimension pieces should give a
category; for example, the corners of a simplex can be assigned the objects of
a category, the edges between the objects can be assigned 1-morphisms, faces
between the edges can be assigned 2-morphisms, and so on. This is a natural
generalization of the data ($Z$) defined in \deref{d:tqft}. In this
viewpoint, the Hilbert space ($Z_N$) and the path integral ($Z_M$) data can
respectively be thought of as objects and morphisms in the category of vector
spaces over complex numbers. Generalizing in this direction requires a notion
of $n$-categories (or a variation of that, called $(\infty, n)$ categories). We
will talk more about higher categories in \chref{c:higher-cat}. 


This chapter was roughly centered around the theme that the subject of TQFT
aims to sidestep the analytical difficulties associated with QFTs. One could
however also argue for the viewpoint that the mathematical study of TQFTs is
actually a step in dealing with analytical difficulties that the current
formulation of QFT suffers from. In either case, we hope the reader is
convinced that the subject of TQFT lies at a very interesting intersection of
modern physics and mathematics.


